\documentclass[11pt]{article}

\usepackage[margin=1in]{geometry}
\usepackage{amsmath, amssymb}
\usepackage{enumitem}

\title{Feature-Based 2D Registration: Algorithmic Description}
\author{}
\date{}

\begin{document}
\maketitle

\section*{1. Inputs}
Two single-channel images are provided as discrete intensity fields, a source image (\texttt{img1}) and a target image (\texttt{img2}), together with a configuration bundle that selects an interest-point method, a local descriptor family, a pairing rule, a motion family, and robust fitting hyperparameters (\texttt{DetectorConfig}, \texttt{DescriptorConfig}, \texttt{MatchConfig}, \texttt{ModelKind}, \texttt{RansacConfig}).

\section*{2. Outputs}
The pipeline returns an estimated $3\times 3$ homogeneous transform mapping source coordinates to target coordinates (\texttt{out.H}), a set of tentative correspondences and an inlier indicator vector supporting the final estimate (\texttt{out.matches}, \texttt{out.inliers}), and a scalar summary of geometric residuals over the inlier set (\texttt{out.mean\_inlier\_error}).

\section*{3. Algorithm Pseudocode}
\begin{enumerate}[label=(\arabic*), leftmargin=*]
  \item Condition the image fields. \\
  \emph{Input:} source and target images. \emph{Output:} smoothed fields and derivatives used by (2), (4), and (5).

  \item Score salient locations. \\
  \emph{Input:} derivatives and/or intensities from (1). \emph{Output:} saliency maps or FAST cues for (3).

  \item Extract spatially separated peaks. \\
  \emph{Input:} saliency from (2). \emph{Output:} ranked keypoints for (4).

  \item Attribute local orientation. \\
  \emph{Input:} keypoints from (3) and derivatives from (1). \emph{Output:} oriented keypoints for (5).

  \item Form local descriptors. \\
  \emph{Input:} oriented keypoints from (4) and fields from (1). \emph{Output:} descriptors and filtered keypoints for (6).

  \item Pair candidates across images. \\
  \emph{Input:} descriptors from (5). \emph{Output:} tentative correspondences for (7).

  \item Estimate motion robustly. \\
  \emph{Input:} correspondences from (6). \emph{Output:} transform and inliers for (8).

  \item Resample and report diagnostics. \\
  \emph{Input:} transform and inliers from (7). \emph{Output:} optional warped image and accuracy summaries.
\end{enumerate}

\section*{Stage (1): Conditioning}
This stage consumes the raw intensity fields (\texttt{img1}, \texttt{img2}) and produces smoothed images and spatial derivatives used by saliency scoring (2), orientation attribution (4), and local description (5).

\begin{equation}
I(y,x)\equiv \text{intensity at row }y\text{ and column }x .
\end{equation}
The discrete field $I$ formalizes pixel access as an evaluation of a sampled scalar function, implemented by the image container and its element accessor (\texttt{operator()}) and used uniformly by filtering, derivative estimation, and interpolation.

\begin{equation}
g_{\sigma}(t)=\frac{1}{\sqrt{2\pi}\sigma}\exp\!\left(-\frac{t^{2}}{2\sigma^{2}}\right),\qquad \sum_{t} g_{\sigma}(t)=1 .
\end{equation}
The one-dimensional Gaussian kernel generated by (\texttt{gaussian\_kernel1d}) is normalized to unit mass, ensuring that subsequent smoothing preserves average intensity while attenuating high-frequency components.

\begin{equation}
(I*g_{\sigma})(x)=\sum_{k=-r}^{r} g_{\sigma}(k)\, I(\operatorname{ref}(x-k)).
\end{equation}
Separable smoothing (\texttt{gaussian\_blur}) is implemented via one-dimensional convolution with reflected boundary handling (\texttt{convolve1d\_reflect}, \texttt{reflect\_index}), producing stable edge behavior without introducing spurious discontinuities.

\begin{equation}
I_{x}=I*S_{x},\qquad I_{y}=I*S_{y}.
\end{equation}
Spatial derivatives (\texttt{sobel\_gradients}) are approximated by Sobel stencils $S_x$ and $S_y$, yielding gradient fields that act as sufficient statistics for second-moment saliency (2) and gradient-based description (5).

\section*{Stage (2): Scoring}
This stage consumes the derivatives and/or intensities from (1) and produces scalar saliency fields used by peak extraction (3).

\begin{equation}
\mathbf{J}(p)=G_{\sigma}*\begin{bmatrix}
I_{x}(p)^{2} & I_{x}(p)I_{y}(p)\\
I_{x}(p)I_{y}(p) & I_{y}(p)^{2}
\end{bmatrix}.
\end{equation}
The local second-moment matrix (\texttt{structure\_tensor}) aggregates gradient outer products under Gaussian smoothing, returning component images that summarize local directional energy in a noise-robust manner.

\begin{equation}
R(p)=\det(\mathbf{J}(p)) - k\,\operatorname{tr}(\mathbf{J}(p))^{2}.
\end{equation}
The Harris response (\texttt{harris\_response}) promotes neighborhoods with strong two-dimensional variation through $\det(\mathbf{J})$ while penalizing anisotropic structure via the trace term, with $k$ controlling the trade-off.

\begin{equation}
\lambda_{\min}(p)=\frac{\operatorname{tr}(\mathbf{J}(p))-\sqrt{\operatorname{tr}(\mathbf{J}(p))^{2}-4\det(\mathbf{J}(p))}}{2}.
\end{equation}
The Shi--Tomasi criterion (\texttt{shitomasi\_response}) uses the smaller eigenvalue of $\mathbf{J}$, directly quantifying the weakest directional gradient energy and thereby emphasizing well-conditioned corners.

\begin{equation}
\exists\ \text{9 contiguous } i \text{ on the 16-point circle: }\ I(p_{i})\ge I(p)+t\ \ \text{or}\ \ I(p_{i})\le I(p)-t .
\end{equation}
The FAST-9 predicate (\texttt{detect\_fast9}) declares a corner when a contiguous arc of length nine on the Bresenham circle is consistently brighter or darker than the center by threshold $t$, enabling efficient rejection and rapid interest-point discovery.

\section*{Stage (3): Peak extraction}
This stage converts saliency outputs from (2) into a ranked, spatially pruned keypoint set consumed by orientation attribution (4).

\begin{equation}
\tau=\alpha \max_{p} R(p).
\end{equation}
A relative cutoff (\texttt{rel\_threshold}) transforms raw saliency into a scale-adaptive threshold, stabilizing the retained keypoint density across images by referencing the maximum response.

\begin{equation}
R(p)\ge \tau\ \wedge\ R(p)\ge \max_{\lVert q-p\rVert_{\infty}\le r} R(q).
\end{equation}
Non-maximum suppression (\texttt{nms\_peaks}) enforces local maximality within a radius-$r$ neighborhood under the $\ell_{\infty}$ metric, ensuring spatial separation of retained keypoints and reducing redundant detections.

\begin{equation}
s(p)=\sum_{j=1}^{16}\left|I(p_{j})-I(p)\right|.
\end{equation}
A contrast-based FAST score used inside (\texttt{detect\_fast9}) provides a ranking signal for suppression and selection by aggregating absolute deviations between the center and circle samples.

\section*{Stage (4): Orientation attribution}
This stage consumes keypoints from (3) and derivatives from (1) to produce oriented keypoints for descriptor formation (5).

\begin{equation}
\bar g_x=\frac{1}{9}\sum_{u=-1}^{1}\sum_{v=-1}^{1} I_x(y+u,x+v),\qquad
\bar g_y=\frac{1}{9}\sum_{u=-1}^{1}\sum_{v=-1}^{1} I_y(y+u,x+v).
\end{equation}
A local gradient mean (\texttt{quick\_orientations}) estimates dominant direction by averaging derivatives in a $3\times 3$ neighborhood, reducing sensitivity to pixel-scale noise while remaining computationally lightweight.

\begin{equation}
\theta=\operatorname{atan2}(\bar g_y,\bar g_x).
\end{equation}
The assigned angle $\theta$ (stored as \texttt{angle}) defines the keypoint orientation, enabling rotation-aligned sampling in (5) and improving repeatability under in-plane rotations.

\section*{Stage (5): Local description}
This stage consumes oriented keypoints from (4) and fields from (1) to produce descriptors and filtered keypoints passed to candidate pairing (6).

\begin{equation}
I(x,y)=(1-a_x)(1-a_y)I_{00}+a_x(1-a_y)I_{10}+(1-a_x)a_yI_{01}+a_xa_yI_{11}.
\end{equation}
Bilinear interpolation (\texttt{bilinear}) provides a continuous sampling model for non-integer coordinates, which is required by rotation-aware descriptors and by inverse-mapped resampling in (8).

\begin{equation}
w(x',y')=\exp\!\left(-\frac{x'^2+y'^2}{2\sigma_w^{2}}\right),\qquad \text{vote}=m(x',y')\,w(x',y').
\end{equation}
The SIFT-like weighting rule in (\texttt{sift128\_descriptor}) down-weights peripheral samples through a Gaussian window while scaling contributions by gradient magnitude $m$, improving stability with respect to small perturbations.

\begin{equation}
\mathbf{d}\leftarrow \frac{\mathbf{d}}{\lVert \mathbf{d}\rVert_2+\varepsilon}.
\end{equation}
Normalization in (\texttt{sift128\_descriptor}) and (\texttt{patch\_descriptor}) enforces scale invariance in descriptor space, reducing sensitivity to global contrast changes and facilitating metric-based matching in (6).

\begin{equation}
d_i \leftarrow \min(\beta, d_i).
\end{equation}
Component-wise clamping in (\texttt{sift128\_descriptor}) limits the influence of dominant orientation bins, which empirically increases robustness to local illumination changes and repeated patterns by preventing a small subset of bins from dominating.

\begin{equation}
\mu=\frac{1}{N}\sum_{k=1}^{N} p_k,\qquad \tilde p_k=p_k-\mu.
\end{equation}
Patch mean-centering in (\texttt{patch\_descriptor}) removes additive brightness bias over a window of $N$ samples, aligning the descriptor with local contrast structure rather than absolute intensity.

\begin{equation}
b_i=\mathbb{1}\{I(x_1,y_1)<I(x_2,y_2)\}.
\end{equation}
BRIEF bit construction (\texttt{brief\_descriptor}) encodes local appearance as binary intensity comparisons, yielding a compact representation suitable for Hamming-distance evaluation in (6).

\begin{equation}
\begin{bmatrix}x'\\y'\end{bmatrix}=
\begin{bmatrix}\cos\theta & -\sin\theta\\ \sin\theta & \cos\theta\end{bmatrix}
\begin{bmatrix}x\\y\end{bmatrix}.
\end{equation}
Rotation of BRIEF test offsets within (\texttt{brief\_descriptor}) aligns the sampling pattern to the orientation from (4), improving repeatability across rotated views without requiring a dense orientation histogram.

\section*{Stage (6): Candidate pairing}
This stage consumes descriptor sets from (5) and produces tentative correspondences passed to robust estimation (7).

\begin{equation}
d(\mathbf{a},\mathbf{b})=\lVert \mathbf{a}-\mathbf{b}\rVert_2.
\end{equation}
The Euclidean metric (\texttt{l2\_dist}) is applied to normalized real-valued descriptors, where normalization renders the distance less sensitive to global scaling and supports nearest-neighbor matching.

\begin{equation}
d(\mathbf{a},\mathbf{b})=\sum_{k}\operatorname{popcount}(a_k\oplus b_k).
\end{equation}
The Hamming metric (\texttt{hamming\_dist}, \texttt{popcount8}) counts differing bits between binary descriptors, supporting efficient pairing for BRIEF representations produced in (5).

\begin{equation}
d_{1}\le \rho\, d_{2}.
\end{equation}
The ratio criterion (\texttt{match\_ratio\_mutual}) accepts a nearest neighbor only when it is sufficiently separated from the second nearest neighbor, reducing ambiguous assignments in repetitive texture and improving the quality of correspondences entering (7).

\section*{Stage (7): Robust motion estimation}
This stage consumes tentative correspondences from (6) and returns a final transform and inlier set used by diagnostics and resampling (8).

\begin{equation}
\mathbf{x}' \sim \mathbf{H}\mathbf{x}.
\end{equation}
Homogeneous mapping (matrix application \texttt{mul}) expresses planar motion in a projective form, with $\mathbf{H}$ representing the selected motion family (\texttt{ModelKind}) and equality defined up to nonzero scale.

\begin{equation}
(x',y')=\left(\frac{u}{w},\frac{v}{w}\right)\quad \text{for}\quad [u,v,w]^{\top}=\mathbf{H}[x,y,1]^{\top}.
\end{equation}
Dehomogenization (\texttt{proj}) converts homogeneous coordinates into Euclidean pixel coordinates, enabling residual evaluation in pixel space and supporting inverse mapping in (8).

\begin{equation}
\mathbf{T}=
\begin{bmatrix}
s&0&-s\mu_x\\
0&s&-s\mu_y\\
0&0&1
\end{bmatrix},
\qquad
s=\frac{\sqrt{2}}{\operatorname{rms}+\varepsilon}.
\end{equation}
Coordinate normalization (\texttt{norm\_matrix\_for\_points}, \texttt{apply\_norm}) improves numerical conditioning by centering points at the origin and scaling them so that the typical radius is $\sqrt{2}$, stabilizing least-squares estimation for affine and projective solvers.

\begin{equation}
\min_{\mathbf{p}}\ \lVert \mathbf{A}\mathbf{p}-\mathbf{b}\rVert_2^2.
\end{equation}
Least-squares fitting (\texttt{least\_squares\_normal}, \texttt{solve\_linear}) estimates model parameters by minimizing the squared residual over an overdetermined linear system, implemented via normal equations and Gaussian elimination with pivoting.

\begin{equation}
\begin{bmatrix}u\\v\end{bmatrix}=
\begin{bmatrix}a&b\\d&e\end{bmatrix}
\begin{bmatrix}x\\y\end{bmatrix}
+
\begin{bmatrix}c\\f\end{bmatrix}.
\end{equation}
The affine motion parameterization used in (\texttt{fit\_affine\_ls}) yields linear constraints in the unknowns $(a,b,c,d,e,f)$ and supports estimation from three or more correspondences after normalization.

\begin{equation}
u=\frac{h_1x+h_2y+h_3}{h_7x+h_8y+1},\qquad
v=\frac{h_4x+h_5y+h_6}{h_7x+h_8y+1}.
\end{equation}
The homography model in (\texttt{fit\_homography\_ls}) represents general planar projective mappings and is required whenever the apparent motion includes perspective components, as in the synthetic demo transform.

\begin{equation}
e_i=\left\lVert \pi(\mathbf{H}\mathbf{x}_i)-\mathbf{y}_i\right\rVert_2.
\end{equation}
Forward transfer error (\texttt{forward\_error}) measures geometric disagreement in the target frame under the estimated transform, providing the residual used for inlier classification in (\texttt{ransac\_fit}).

\begin{equation}
e_i=\left\lVert \pi(\mathbf{H}\mathbf{x}_i)-\mathbf{y}_i\right\rVert_2+\left\lVert \pi(\mathbf{H}^{-1}\mathbf{y}_i)-\mathbf{x}_i\right\rVert_2.
\end{equation}
Symmetric transfer error (\texttt{symmetric\_error}, \texttt{invert}) penalizes mismatch in both mapping directions, which can improve stability when residuals are directionally biased by noise or local distortions.

\begin{equation}
\text{cost}=\sum_i \min(e_i^2,\tau^2).
\end{equation}
Robust scoring in (\texttt{ransac\_fit}) uses a truncated quadratic objective that limits the leverage of large residuals while still favoring hypotheses with both broad consensus and low inlier error, after which the model is refit using the final inlier set.

\section*{Stage (8): Resampling and diagnostics}
This stage consumes the transform and inliers from (7) and produces optional warped imagery and summary accuracy statistics.

\begin{equation}
\mathbf{p}=\pi(\mathbf{H}^{-1}\mathbf{q}),\qquad I_{\text{out}}(\mathbf{q})=I_{\text{src}}(\mathbf{p}).
\end{equation}
Inverse-mapped warping (\texttt{warp\_perspective}) assigns each output pixel $\mathbf{q}$ by sampling the source at the inverse-projected location $\mathbf{p}$, yielding dense coverage of the output grid and leveraging bilinear sampling from (5).

\begin{equation}
\bar e=\frac{1}{|\mathcal{I}|}\sum_{i\in\mathcal{I}} e_i.
\end{equation}
The reported mean inlier error computed in the registration wrapper (\texttt{register\_pair}) averages the residuals over the inlier set $\mathcal{I}$ obtained in (7), producing a compact measure of geometric consistency.

\begin{equation}
\text{corner\_err}=\frac{1}{4}\sum_{j=1}^{4}\left\lVert \pi(\mathbf{H}\mathbf{c}_j)-\pi(\mathbf{H}_{\text{gt}}\mathbf{c}_j)\right\rVert_2.
\end{equation}
The corner-based evaluation statistic (\texttt{corner\_error}) compares the estimated transform to a ground-truth transform $\mathbf{H}_{\text{gt}}$ on the set of image corners $\{\mathbf{c}_j\}$, yielding an interpretable pixel-scale discrepancy that reflects global alignment quality.

\end{document}
